% ===== PRÉAMBULE CANONIQUE — à coller VERBATIM en tête de CHAQUE .tex =====
\documentclass[11pt,a4paper]{article}
\usepackage[utf8]{inputenc}\usepackage[T1]{fontenc}\usepackage{lmodern}
\usepackage[french]{babel}
\usepackage{amsmath,amssymb,mathtools}\usepackage{icomma}
\usepackage{siunitx}\sisetup{locale=FR}  % \qty{}{}, \si{}, \ang{}
\usepackage{xcolor}\usepackage{colortbl}
\usepackage{tikz}\usepackage{pgfplots}\pgfplotsset{compat=1.18}
\usepackage[siunitx,europeanresistors]{circuitikz} % circuits (élec.)
\usepackage[most]{tcolorbox}\usepackage{enumitem}\usepackage{array,booktabs}
\usepackage[a4paper,margin=2.2cm]{geometry}
\usetikzlibrary{arrows.meta,calc,angles,quotes,positioning,patterns,%
  backgrounds,shapes.geometric,decorations.pathmorphing,decorations.markings,intersections}
\definecolor{primary}{RGB}{222,49,99}\definecolor{primarydark}{RGB}{150,22,55}
\definecolor{defcol}{RGB}{0,114,178}\definecolor{methcol}{RGB}{52,92,148}
\definecolor{excol}{RGB}{0,135,90}\definecolor{attncol}{RGB}{200,40,55}
\tcbset{boxbase/.style={enhanced,breakable,boxrule=0.9pt,arc=2.5pt,
  left=3mm,right=3mm,top=2.3mm,bottom=2.3mm,fonttitle=\bfseries,coltitle=white}}
% --- boîtes TUTORIEL (noms EXACTS, ne pas renommer) ---
\newtcolorbox{cle}[1][]{boxbase,colback=defcol!6,colframe=defcol,
  title={\textsc{À retenir}\ifx\relax#1\relax\relax\else\textmd{ : \itshape #1}\fi}}
\newtcolorbox{methode}[1][]{boxbase,colback=methcol!7,colframe=methcol,sharp corners,
  title={\textsc{Méthode}\ifx\relax#1\relax\relax\else\textmd{ --- \itshape #1}\fi}}
\newtcolorbox{exemplebox}[1][]{boxbase,colback=excol!5,colframe=excol,
  title={\textsc{Exemple}\ifx\relax#1\relax\relax\else\textmd{ : \itshape #1}\fi}}
\newtcolorbox{piege}[1][]{boxbase,colback=attncol!6,colframe=attncol,
  title={\textsc{Piège}\ifx\relax#1\relax\relax\else\textmd{ : \itshape #1}\fi}}
\newtcolorbox{remarque}[1][]{boxbase,colback=black!4,colframe=black!45,
  title={\textsc{Remarque}\ifx\relax#1\relax\relax\else\textmd{ : \itshape #1}\fi}}
% --- boîtes EXERCICE / CORRIGÉ (noms EXACTS, auto-comptées) ---
\newtcolorbox[auto counter]{exo}[1][]{boxbase,colback=primary!4,colframe=primary,
  title={\textsc{Exercice}~\thetcbcounter\ifx\relax#1\relax\relax\else\textmd{ --- \itshape #1}\fi}}
\newtcolorbox[auto counter]{corrige}[1][]{boxbase,colback=excol!4,colframe=excol,
  title={\textsc{Corrigé}~\thetcbcounter\ifx\relax#1\relax\relax\else\textmd{ --- \itshape #1}\fi}}
\newcommand{\vv}[1]{\vec{#1}}\newcommand{\ds}{\displaystyle}
\newcommand{\R}{\mathbb{R}}\newcommand{\N}{\mathbb{N}}\newcommand{\Z}{\mathbb{Z}}
\begin{document}
\sisetup{per-mode=symbol}

\begin{corrige}[le projecteur de diapositives]
\begin{enumerate}[label=\alph*)]
  \item $|\gamma|=\dfrac{|i|}{o}=\dfrac{120}{2,4}=50$ ; image réelle renversée $\Rightarrow \gamma=-50$.
        De $\gamma=-v/u$ : $v=-\gamma\,u=50\times 6=\qty{300}{\cm}=\qty{3}{\m}$ (l'écran est à \qty{3}{\m}).
  \item $\dfrac{1}{f}=\dfrac{1}{u}+\dfrac{1}{v}=\dfrac{1}{6}+\dfrac{1}{300}=\dfrac{50+1}{300}=\dfrac{51}{300}
        \Rightarrow f=\dfrac{300}{51}\approx\qty{5.9}{\cm}$. On a bien $f<u<2f$ (projecteur).
\end{enumerate}
\end{corrige}

\begin{corrige}[l'appareil photo saisit un immeuble]
\begin{enumerate}[label=\alph*)]
  \item $\dfrac{1}{v}=\dfrac{1}{f}-\dfrac{1}{u}=\dfrac{1}{0,05}-\dfrac{1}{50}=20-0{,}02=19{,}98
        \Rightarrow v=\qty{0.0500}{\m}\approx\qty{50.1}{\mm}$ : pour un objet lointain, $v\approx f$.
  \item $\gamma=-\dfrac{v}{u}=-\dfrac{0,05005}{50}\approx-1{,}0\times10^{-3}$ ;
        $i=\gamma\,o=-10^{-3}\times 20\approx\qty{-20}{\mm}$. Image réelle, renversée, très réduite.
\end{enumerate}
\end{corrige}

\begin{corrige}[grossissement d'un microscope]
\begin{enumerate}[label=\alph*)]
  \item $|i_1|=|\gamma_1|\times o=40\times\qty{0.01}{\mm}=\qty{0.4}{\mm}$ (réelle, renversée, agrandie).
  \item $G\approx|\gamma_1|\times G_{\text{oculaire}}=40\times 10=400$ : l'objet est perçu $\approx 400$
        fois plus grand.
\end{enumerate}
\end{corrige}

\begin{corrige}[la loupe en action]
\begin{enumerate}[label=\alph*)]
  \item $\dfrac{1}{v}=\dfrac{1}{5}-\dfrac{1}{4}=\dfrac{4-5}{20}=-\dfrac{1}{20}\Rightarrow v=-\qty{20}{\cm}$ :
        $v<0$, image \textbf{virtuelle}.
  \item $\gamma=-\dfrac{-20}{4}=+5$ ; $i=\gamma\,o=5\times 3=\qty{15}{\mm}$. Image droite, agrandie
        $5\times$ : la loupe grossit.
\end{enumerate}
\end{corrige}

\begin{corrige}[faire la mise au point]
D'après $\dfrac{1}{v}=\dfrac{1}{f}-\dfrac{1}{u}$ :
\begin{enumerate}[label=\alph*)]
  \item quand $u$ diminue, $\dfrac{1}{u}$ augmente, donc $\dfrac{1}{v}$ diminue et \textbf{$v$ augmente}
        : l'image se forme plus loin derrière l'objectif.
  \item il faut \textbf{éloigner l'objectif du capteur} (l'objectif « sort » de l'appareil pour un
        objet proche ; pour un objet à l'infini, $v=f$, objectif au plus près).
\end{enumerate}
\end{corrige}

\end{document}
